Un bassin destiné à l'élevage des truites contient \np{30000} litres d'eau, régulièrement renouvelée.

	Pour vérifier la présence éventuelle de substances polluantes dans le bassin, on procède à des analyses régulières.

	On étudie dans les deux parties ci-dessous un cas de contamination de l'eau par une substance polluante à l'aide de deux modèles différents.

	Les deux parties sont indépendantes.

	\subsection*{\centering Partie A : modèle discret}

		On note $V_{n}$ le volume, en litre, de la substance polluante présente dans le bassin, $n$ heures après le début du phénomène de pollution.

		On a ainsi \quad $V_{0}=0$.

		On admet que pour tout entier naturel $n$,\quad $V_{n+1}=0,995 V_{n}+6$.

		\begin{enumerate}
			\item Calculer $V_{1}$ et $V_{2}$.

			\item \begin{minipage}[t]{0.55\linewidth}
				Recopier et compléter le programme ci-contre pour qu'il renvoie le volume en litre de substance polluante présente dans le bassin après $n$ heures.
				\item Montrer par récurrence que pour tout entier naturel $n$,  $V_{n} \leqslant  V_{n+1} \leqslant  1200$.

			\end{minipage}\hfill
			\begin{tblr}[T]{colspec={X[font=\ttfamily,5cm]},vlines, hline{1,Z}={solid}, rowsep=0pt}
				def volume(n):\\
				\quad v = ...\\
				\quad for k in range(n):\\
				\quad \quad	v = ...\\
				\quad return v\\
			\end{tblr}


			\item Justifier que la suite $(V_{n})$ converge, et calculer sa limite.
		\end{enumerate}

	\subsection*{\centering Partie B : modèle continu}

		On note $v(t)$ le volume, en litre, de la substance polluante présente dans le bassin à l'instant $t$, en heure.

		On admet que la fonction $v$ est définie et dérivable sur l'intervalle $[0 ;+\infty[$, et on note $v'$ sa fonction dérivée.

		On considère que le phénomène de pollution débute au temps $t=0$; ainsi $v(0)=0$.

		On admet que la fonction $v$ est solution de l'équation différentielle :
		\[(E) : \quad y'=-0,005 y+6,\]
		où $y$ est une fonction inconnue et où $y'$ est la fonction dérivée de $y$.

		\begin{enumerate}
			\item \begin{enumerate}
				\item Déterminer les solutions de l'équation différentielle $(E)$ sur $[0 ;+\infty[$.

				\item Montrer que pour tout réel $t$ appartenant à l'intervalle $[0 ;+\infty[$ :
				\[v(t)=1200\left(1-\e^{-0,005 t}\right).\]

				\item Calculer la limite de $v$ en $+\infty$.

				\item Déterminer le sens de variation de la fonction $v$ sur l'intervalle $[0 ;+\infty[$.
			\end {enumerate}

			\item Un taux de substance polluante supérieur à $5\,\%$ du volume du bassin nécessite un nettoyage complet.
			Selon ce modèle, le propriétaire devra-t-il procéder à ce nettoyage?

			\item Pour pouvoir dépolluer rapidement le bassin et éviter de lourdes conséquences sur l'élevage, il est préconisé que le volume de la substance polluante ne dépasse pas 50 litres.

			Déterminer la valeur exacte de l'instant au-delà duquel le volume de la substance dépasse 50 litres dans le bassin.

			En donner la valeur en heure et minute, arrondie à la minute.
		\end{enumerate}
